\documentclass{article}
\usepackage[T1]{fontenc}
\usepackage{lmodern}
\usepackage{graphicx}
\usepackage{amsmath,amssymb}
\usepackage[a4paper,margin=2cm]{geometry}
\usepackage[absolute,overlay]{textpos}
\setlength{\TPHorizModule}{1mm}
\setlength{\TPVertModule}{1mm}
\usepackage[indonesian]{babel}
\usepackage{hyperref}
\setlength{\emergencystretch}{2em}
% unit: Halaman 46

\begin{document}

% === Halaman 46 (python/native) ===

Counter Mode

• Mode counter tidak melakukan perantaian (chaining) seperti pada CBC. 

• Counter adalah sebuah nilai berupa blok bit yang ukurannya sama dengan 

ukuran blok plainteks. 

• Nilai counter harus berbeda dari setiap blok yang dienkripsi. Pada mulanya, 

untuk enkripsi blok pertama, counter diinisialisasi dengan sebuah nilai. 

• Selanjutnya, untuk enkripsi blok-blok berikutnya counter dinaikkan 

(increment) nilainya satu  (counter  counter + 1). 

46

\end{document}
